\chapter{Unit 5: Isometric Views \& 3D Modeling Fundamentals}

\section{A. What You'll Learn}
Isometric projection is a pictorial drawing method enabling a 3D object to be visualized on a 2D plane surface. In this unit, you will master:
\begin{itemize}
    \item Isometric principles, Isometric Axes ($120^\circ$ separation, $30^\circ$ receding axes), Isometric Lines, and Non-Isometric Lines.
    \item Mathematical derivation of the **Isometric Scale Reduction Factor ($0.816$)**.
    \item Differentiation between **Isometric View (Drawing)** (True Scale $1:1$) and **Isometric Projection** (Reduced Scale $0.816$).
    \item The **Box Method** for constructing prisms, pyramids, and composite/stacked solids.
    \item Isometric dimensioning rules and obliqued text angles ($-30^\circ / +30^\circ$).
    \item AutoCAD 3D Modeling Coordinate Systems (`3P UCS`), 3D Solid Primitives (`BOX`, `CYLINDER`, `CONE`, `SPHERE`), Extrude (`EXTRUDE`), Revolve (`REVOLVE`), and Presspull (`PRESSPULL`).
\end{itemize}

---

\section{B. Isometric Principles \& Terminology}

\begin{definition}[Isometric Projection]
A type of axonometric pictorial projection where three mutually perpendicular edges of a cube are equally inclined to the plane of projection. Consequently, all three axes are foreshortened by the exact same amount.
\end{definition}

\subsection{Key Isometric Characteristics}
\begin{itemize}
    \item \textbf{Isometric Axes:} Three principal lines meeting at a point and making an angle of **$120^\circ$** with each other.
    \begin{itemize}
        \item Vertical Axis: Represents height.
        \item Right Axis ($30^\circ$ to horizontal): Represents length/width.
        \item Left Axis ($30^\circ$ to horizontal): Represents width/length.
    \end{itemize}
    \item \textbf{Isometric Lines:} Any line parallel to an isometric axis. Dimensions along isometric lines can be measured directly.
    \item \textbf{Non-Isometric Lines:} Any line NOT parallel to an isometric axis (e.g., slant edges of pyramids, diagonals of cubes). **Non-isometric lines cannot be measured directly**; their endpoints must be located using isometric coordinates.
    \item \textbf{Isometric Planes:} The three faces of a cube formed by pairs of isometric axes (Top plane, Left vertical plane, Right vertical plane).
\end{itemize}

---

\section{C. Isometric Scale Derivation \& Scale Factor}

In a true isometric projection, edges are foreshortened because the cube is tilted relative to the projection plane.

\begin{keyequation}[Mathematical Derivation of Isometric Scale Factor]
Consider a cube top face tilted at $45^\circ$ to the true horizontal line and $30^\circ$ to the isometric axis line:
\begin{equation}
    \frac{\text{Isometric Length}}{\text{True Length}} = \frac{\cos 45^\circ}{\cos 30^\circ} = \frac{1/\sqrt{2}}{\sqrt{3}/2} = \sqrt{\frac{2}{3}} \approx \mathbf{0.8165}
\end{equation}

\begin{equation}
    \mathbf{\text{Isometric Length} = 0.816 \times \text{True Length}}
\end{equation}
\end{keyequation}

\subsection{Isometric View vs. Isometric Projection}

\begin{table}[htbp]
\centering
\small
\begin{tabularx}{\textwidth}{l X X}
\toprule
\textbf{Parameter} & \textbf{Isometric View (Drawing)} & \textbf{Isometric Projection} \\
\midrule
\textbf{Scale Used} & **True Scale** ($1:1$, actual physical dimensions) & **Isometric Scale** ($0.816 \times \text{True Length}$) \\
\textbf{Size} & Larger ($\approx 22.5\%$ larger than projection) & Real foreshortened size \\
\textbf{Constructibility} & Fast and simple to construct directly & Requires constructing an isometric scale line \\
\textbf{Standard Usage} & Standard practice in AutoCAD and CAD drafting & Theoretical descriptive geometry examinations \\
\bottomrule
\end{tabularx}
\caption{Comparison of Isometric View and Isometric Projection.}
\label{tab:iso-view-vs-proj}
\end{table}

---

\section{D. Construction Methods for Isometric Solids}

\subsection{The Box Method (Enclosing Box Method)}
Solids are constructed by imagining the object enclosed inside a rectangular box of overall Length ($L$), Width ($W$), and Height ($H$).

\subsubsection{Procedure for Prisms}
1. Draw the three isometric axes meeting at origin ($120^\circ$ apart, two receding at $30^\circ$).
2. Construct the base polygon on the isometric plane inside an enclosing rectangle.
3. Erect vertical isometric lines from all base corners equal to height $H$.
4. Connect top points to form top face.

\subsubsection{Procedure for Pyramids}
1. Enclose base polygon in a rectangle on the bottom isometric plane.
2. Draw diagonals across the base to locate center point $O$.
3. From $O$, erect a vertical line equal to axis height $H$ to locate Apex $P$.
4. Connect Apex $P$ to base corners. (Note: Slant edges are **non-isometric lines**).

\subsection{Composite / Stacked Solids}
For objects stacked on top of one another (e.g., a cylinder on a cube):
1. Construct the bottom solid completely.
2. Locate center point of top face of bottom solid using diagonals.
3. Construct base of upper solid centered on that point.
4. Draw upper solid and erase hidden obscured lines.

---

\section{E. AutoCAD 3D Modeling Commands}

Transitioning from 2D drafting to 3D solid modeling requires mastering coordinate systems and 3D primitives.

\begin{autocadcmd}[3P UCS (3-Point User Coordinate System)]
Moves the coordinate system working plane to any face of a 3D object.
\begin{itemize}
    \item \textbf{Command:} \texttt{UCS} $\to$ Type \texttt{3} (3 Point)
    \item \textbf{Workflow:} Specify new Origin ($0,0,0$) $\to$ Specify point on positive X-axis $\to$ Specify point on positive Y-axis.
\end{itemize}
\end{autocadcmd}

\subsection{3D Solid Primitives}
\begin{itemize}
    \item \texttt{BOX}: Creates a rectangular solid prism ($L, W, H$).
    \item \texttt{CYLINDER}: Creates a cylinder with defined base radius and height.
    \item \texttt{CONE}: Creates a cone with defined base radius and height.
    \item \texttt{SPHERE}: Creates a solid ball with defined radius.
\end{itemize}

\begin{autocadcmd}[EXTRUDE]
Converts closed 2D polylines, circles, or regions into 3D solids by stretching along Z-axis.
\begin{itemize}
    \item \textbf{Command:} \texttt{EXTRUDE} or \texttt{EXT}
    \item \textbf{Input:} Closed polyline, circle, or region.
    \item \textbf{Workflow:} Select object $\to$ Enter height $\to$ Enter taper angle (optional).
\end{itemize}
\end{autocadcmd}

\begin{autocadcmd}[REVOLVE]
Creates a solid by sweeping a 2D closed profile around an axis line.
\begin{itemize}
    \item \textbf{Command:} \texttt{REVOLVE} or \texttt{REV}
    \item \textbf{Workflow:} Select 2D profile $\to$ Select axis line $\to$ Enter angle of revolution (usually $360^\circ$).
\end{itemize}
\end{autocadcmd}

\begin{autocadcmd}[PRESSPULL]
Dynamically detects enclosed coplanar boundary areas and adds/subtracts 3D volume.
\begin{itemize}
    \item \textbf{Command:} \texttt{PRESSPULL}
    \item \textbf{Advantage over Extrude:} Does not require a polyline; automatically detects enclosed areas bounded by coplanar lines.
\end{itemize}
\end{autocadcmd}

---

\section{F. Chapter 5 Practice Problems \& MCQs}

\subsection{Practice Questions}

\begin{vivaqa}[Question 1 (Basic)]
What are non-isometric lines? Can dimensions be measured directly along them?\\[4pt]
\textbf{Answer:} Non-isometric lines are lines that are NOT parallel to any of the three isometric axes (e.g., slant edges of a pyramid). Dimensions **cannot** be measured directly along them; endpoints must be plotted using isometric coordinates.
\end{vivaqa}

\begin{vivaqa}[Question 2 (Moderate)]
Derive the isometric scale factor ($0.816$).\\[4pt]
\textbf{Answer:} Ratio of isometric length to true length $= \frac{\cos 45^\circ}{\cos 30^\circ} = \frac{1/\sqrt{2}}{\sqrt{3}/2} = \sqrt{\frac{2}{3}} \approx \mathbf{0.816}$.
\end{vivaqa}

\begin{vivaqa}[Question 3 (Exam Level)]
Calculate the isometric length corresponding to a true length of $90\text{ mm}$.\\[4pt]
\textbf{Answer:}  
$$\text{Isometric Length} = 0.816 \times \text{True Length} = 0.816 \times 90\text{ mm} = \mathbf{73.44\text{ mm}}$$
\end{vivaqa}

\subsection{Multiple-Choice Questions (MCQs)}

1. What is the angle between any two principal isometric axes?  
   A. $30^\circ$  
   B. $60^\circ$  
   C. $90^\circ$  
   D. $120^\circ$  
   \textbf{Answer: D} — Isometric axes meet at $120^\circ$ intervals.

2. The ratio of isometric length to true length is approximately:  
   A. $0.577$  
   B. $0.707$  
   C. $0.816$  
   D. $0.866$  
   \textbf{Answer: C} — The ratio is $\sqrt{2/3} \approx 0.816$.

3. In AutoCAD, an Isometric View (Drawing) is created using:  
   A. Reduced Isometric Scale  
   B. True Scale ($1:1$)  
   C. Double Scale  
   D. Architectural Scale  
   \textbf{Answer: B} — Isometric Views use True Scale ($1:1$).

4. Which 3D modeling command sweeps a closed 2D profile around an axis line?  
   A. \texttt{EXTRUDE}  
   B. \texttt{REVOLVE}  
   C. \texttt{PRESSPULL}  
   D. \texttt{SWEEP}  
   \textbf{Answer: B} — \texttt{REVOLVE} sweeps profiles around an axis.

5. What receding angle do horizontal isometric axes make with the true horizontal reference line?  
   A. $15^\circ$  
   B. $30^\circ$  
   C. $45^\circ$  
   D. $60^\circ$  
   \textbf{Answer: B} — Receding isometric axes make a $30^\circ$ angle with horizontal.

---

\section{G. Unit 5 Quick Revision Checklist}
\begin{revisionchecklist}
\begin{itemize}
    \item $\text{Isometric Axes Angle} = 120^\circ \text{ between axes}, \quad 30^\circ \text{ to horizontal}$.
    \item $\text{Isometric Scale Ratio} = \sqrt{2/3} = \frac{\cos 45^\circ}{\cos 30^\circ} \approx 0.816$.
    \item $\text{Isometric View} \to \text{True Scale (1:1)}$; $\text{Isometric Projection} \to \text{Reduced Scale (0.816)}$.
    \item $\text{Non-Isometric Lines} \to \text{Cannot be measured directly; locate endpoints via box coordinates}$.
    \item $\text{Box Method} \to \text{Enclose solid in L, W, H rectangular box}$.
    \item $\text{AutoCAD 3D: } \texttt{EXTRUDE (linear stretch)}, \quad \texttt{REVOLVE (axis sweep)}, \quad \texttt{PRESSPULL (boundary detection)}$.
\end{itemize}
\end{revisionchecklist}
